\documentclass[12pt]{article}

\usepackage[utf8x]{inputenc}
\usepackage[spanish]{babel}

\usepackage{amssymb,amsmath,amsthm,amsfonts}
\usepackage{calc}
\usepackage{graphicx}
\usepackage{subfigure}
\usepackage{gensymb}
\usepackage{natbib}
\usepackage{url}
\usepackage[utf8x]{inputenc}
\usepackage{amsmath}
\usepackage{graphicx}
\graphicspath{{images/}}
\usepackage{parskip}
\usepackage{fancyhdr}
\usepackage{vmargin}
\setmarginsrb{3 cm}{2.5 cm}{3 cm}{2.5 cm}{1 cm}{1.5 cm}{1 cm}{1.5 cm}

\title{Tu titulo}					% Titulo
\author{Tu nombre}					% Autor
\date{\today}						% Fecha


\makeatletter
\let\thetitle\@title
\let\theauthor\@author
\let\thedate\@date
\makeatother

\pagestyle{fancy}
\fancyhf{}
\rhead{\theauthor}
\lhead{\thetitle}
\cfoot{\thepage}

\begin{document}

%%%%%%%%%%%%%%%%%%%%%%%%%%%%%%%%%%%%%%%%%%%%%%%%%%%%%%%%%%%%%%%%%%%%%%%%%%%%%%%%%%%%%%%%%

\begin{titlepage}
	\centering
    \vspace*{0.0 cm}
    \includegraphics[scale = 0.13]{Logo_Uchile_modern.png}\\[1.0 cm]	% Logo Universidad
    \textsc{\LARGE Universidad de Chile}\\[2.0 cm]	% Nombre Universidad
	\textsc{\Large CC3001-02}\\[0.5 cm]				% Codigo Curso
	\textsc{\large Algoritmos y Estructuras de Datos}\\[0.5 cm]		% Nombre Curso
	\rule{\linewidth}{0.2 mm} \\[0.4 cm]
	{ \huge \bfseries \thetitle}\\
	\rule{\linewidth}{0.2 mm} \\[1.5 cm]
	
	\begin{minipage}{0.4\textwidth}
		\begin{center} \large
			\emph{Autor:}\\
			\theauthor\linebreak
			\end{center}
	\end{minipage}\\[2 cm]
	
	{\large \thedate}\\[2 cm]
 
	\vfill
	
\end{titlepage}

%%%%%%%%%%%%%%%%%%%%%%%%%%%%%%%%%%%%%%%%%%%%%%%%%%%%%%%%%%%%%%%%%%%%%%%%%%%%%%%%%%%%%%%%%

\tableofcontents
\pagebreak

%%%%%%%%%%%%%%%%%%%%%%%%%%%%%%%%%%%%%%%%%%%%%%%%%%%%%%%%%%%%%%%%%%%%%%%%%%%%%%%%%%%%%%%%%

\section{Introducción}
\section{Nudo}
\section{blabla}
\section{más apartados}
\section{Introduction}
There is a theory which states that if ever anyone discovers exactly what the Universe is for and why it is here, it will instantly disappear and be replaced by something even more bizarre and inexplicable.
There is another theory which states that this has already happened.
\section{más apartados}
\section{Lo que sea}
There is a theory which states that if ever anyone discovers exactly what the Universe is for and why it is here, it will instantly disappear and be replaced by something even more bizarre and inexplicable.
There is another theory which states that this has already happened.
\section{más apartados}
PruebaPruebaPruebaPruebaPrueba.
\section{más apartados}
\section{explicacion}
\section{más apartados}
\section{otro apartado}

Definición.

Un helióstato es un conjunto de espejos que establecen una superficie grande y se mueven sobre uno o dos ejes, normalmente en montura altacimutal, lo que permite, con los movimientos apropiados, mantener el reflejo de los rayos solares que inciden sobre él, se fijen en todo momento en un punto, o diminuta superficie descomponiendo en el rayo reflejado el movimiento diurno terrestre, teniendo como objetivo seguir el movimiento del sol. Haciendo esto, los rayos que refleja el helióstato pueden ser dirigidos hacia un solo punto durante todo el día.


Utilización.

Se utilizan fundamentalmente en observaciones astronómicas para mantener fija la imagen del Sol o de un astro sobre el aparato de observación. En este caso suelen ser de pequeñas dimensiones. También se utilizan en centrales solares termoeléctricas para concentrar la energía solar sobre el receptor, y conseguir así altas temperaturas. Estos helióstatos suelen ser grandes, llegando a tener más de 120 m2.

En experimentación y pruebas de materiales a altas temperaturas, un conjunto suficientemente alto de helióstatos puede concentrar los rayos solares hasta conseguir temperaturas de más de 2000 ºC.


Desarrollo de los primeros helióstatos.

Los primeros helióstatos considerados como elementos industriales se desarrollaron a los inicios de la década de los ochenta para las plantas experimentales termosolares de receptor central, con el propósito de probar la viabilidad de la energía solar térmica en los procesos de producción de electricidad a escala industrial.


Helióstato con sensor de reflexión.

[5]

El helióstato con sensor de reflexión es un helióstato perteneciente a un campo solar que refleja los haces de luz que llegan a él dotado de un mecanismo de seguimiento solar. Se trata de una invención que pertenece dentro del área de la termotecnia, al campo de la producción de energía a partir de la radiación solar.

[5]

Desde mediados del siglo 20 se vienen realizando investigaciones para intentar transformar esa energía en electricidad. Es por esto que se han desarrollado placas fotovoltaicas que producen directamente electricidad cuando su superficie es convenientemente activada por la luz, y distintos tipos de colectores de calor que concentrando haces de luz sobre una tubería o sobre un receptor central, que contiene un fluido logran alcanzar temperaturas suficientes como para producir grandes cantidades de vapor de agua que genera electricidad a través de una turbina, normalmente en un ciclo de Rankine.

Dada la baja potencia específica por unidad de superficie de la radiación solar, para aprovechar esta energía se debe concentrar un gran número de haces de luz sobre un mismo punto, lo que tradicionalmente se realiza por medio de espejos orientados sobre un depósito o sobre una tubería a modo de colector. En este caso, la radiación es por concentración indirecta, ya que para alcanzar su objetivo los rayos han de rebotar previamente en el espejo. [5]

[1] [2]
Funcionamiento de la central solar

Una central solar de tipo torre central, está formada por un campo de helióstatos o espejos direccionales de grandes dimensiones que reflejan la luz del sol y concentran los haces reflejados en una caldera situada sobre una torre de gran altura.



En la caldera, el aporte calorífico obtenido es mezclado con un fluido térmico. Dicho fluido es conducido hacia un generador de vapor donde transfiere el calor a otro fluido. Este se convierte en vapor, y acciona los álabes del grupo turbina-alternador generando energía eléctrica. El fluido es posteriormente condensado en un aerocondensador para repetir el ciclo.

La producción de una central térmica depende de una serie de factores:

La cantidad de horas que este fluido esté expuesto al sol.
El lugar donde la fábrica esté situada.
La calidad de los depósitos de almacenamiento del fluido.

La energía producida, después de ser transformada, es transportada mediante líneas a la red general. [8]




Energía Solar Térmica — “TORRE CENTRAL”
La torre central está compuesta de un campo de helióstatos los cuales siguen automáticamente al sol.
Sus costos de instalación son superiores a los de una central térmica convencional, es por eso que no se ha comercializado.
La potencia de la torre central va de los 10 a los 200MW, pero hasta ahora su utilización ha sido solamente de investigación.
Su buen funcionamiento para obtener el máximo aprovechamiento depende de la latitud y climatología (situación geográfica) a la que está sometida.
La configuración de las torres son de dos tipos:
Una en la que los helióstatos rodean completamente a la torre central.
En la otra, los helióstatos están colocados al norte de la torre.
BENEFICIOS DE LA TORRE CENTRAL:
Conforme se haya hecho la inversión y la instalación, ya no generará gastos mayores, el único será el del mantenimiento.
Energía eléctrica prácticamente gratis.
No requiere el uso de combustibles, lo que evita el riesgo de almacenar combustibles.
No contamina, ya que la energía solar no produce desechos, residuos ni vapores.
Menos ruidosa que una termo convencional.
INCONVENIENTES:
Se debe de instalar en lugares en donde la radiación del sol sea durante el mayor tiempo posible durante todo el día.
Su rendimiento es menor que el de otros sistemas, en comparación con el de otros sistemas. Su eficiencia es de apenas de un 16 a un 20 por ciento.
Todo su sistema mecánico es más complejo que otros sistemas.
Se necesitan acumuladores de calor para poder usar este dispositivo para cuando no haya radiación.
El funcionamiento de la torre es similar al de una termoeléctrica, ya que en la torre se coloca una caldera, y dentro de esta se coloca algún líquido (agua, aceites, etc.) y este al recibir la radiación del sol eleva su temperatura de 300 a 1500 grados centígrados y esto provoca la evaporación del líquido, de ahí pasa a un acumulador, en donde se almacena energía térmica, después pasa a la turbina.
En la turbina se genera la energía eléctrica, y de ahí pasa a un generador y un transformador y finalmente a una subestación; mientras que el vapor de agua pasa a un condensador y a un sistema de enfriamiento para posteriormente pasar por una bomba a la caldera donde el ciclo se repite. [10] [11] [12]


 
Paneles Fotovoltaicos

Los paneles fotovoltaicos o módulos fotovoltaicos, también llamados a veces paneles solares, están formados por un conjunto de celdas o células fotovoltaicas que producen electricidad a partir de la luz o radiación solar que incide sobre ellos.
Los paneles fotovoltaicos consisten en una red de células conectadas como circuito en serie para aumentar la tensión de salida hasta el valor deseado (usualmente se utilizan 12V ó 24V) a la vez que se conectan varias redes como circuito paralelo para aumentar la corriente eléctrica que es capaz de proporcionar el dispositivo. Poseen una vida útil de 25 a 30 años.
El tipo de electricidad que proporcionan los paneles fotovoltaicos es de corriente continua.
La primera generación de paneles fotovoltaicos tenían una gran superficie de cristal simple capaz de generar energía eléctrica a partir de fuentes de luz con longitudes de onda similares a las que llegan a la superficie de la Tierra provenientes del Sol. Sus células fotovoltaicas se fabrican usualmente utilizando un proceso de difusión con obleas de silicio.
Un módulo o panel fotovoltaico funciona por el efecto fotoeléctrico. Este tipo de paneles fotovoltaicos se suelen montar agrupándolos en solares o parques fotovoltaicos y su efectividad depende tanto de su orientación hacia el sol como de su inclinación con respecto a la horizontal. Para ahorrar gastos de instalación y mantenimiento, los paneles fotovoltaicos se suelen montar con orientación e inclinación fija, tratando de optimizarlos al máximo en función de la latitud.
Los paneles fotovoltaicos se destinan por ejemplo a las siguientes aplicaciones:
Centrales conectadas a red con subvención a la producción
Estaciones repetidoras de radio.
Electrificación de pueblos en áreas remotas
Instalaciones médicas en áreas rurales.
Corriente eléctrica para casas de campo.
Sistemas de comunicaciones de emergencia.
Sistemas de vigilancia de datos ambientales y de calidad del agua.
Faros, boyas y balizas de navegación marítima.
Bombeo para sistemas de riego
Vehículos de recreo.
Sistemas para cargar acumuladores de barcos.
Fuente de energía para naves espaciales.
Postes SOS (teléfonos de emergencia de carretera). [27]



PSA.

El Horno Solar es una de las instalaciones que forma parte de la Plataforma Solar de Almería. Se compone de un helióstato plano que realiza un seguimiento solar continuo, un atenuador o persiana, un espejo parabólico concentrador y la zona de ensayos situada en el foco del concentrador.
            	El helióstato refleja los rayos solares incidentes sobre el disco parabólico, éste a su vez los refleja sobre su foco, donde se encuentra situada el área de ensayos. Un atenuador situado entre el helióstato y el concentrador regula la cantidad de luz. La mesa de ensayos se encuentra bajo el foco y permite su movimiento en todas direcciones para posicionar las probetas con precisión.
            	Los helióstatos están formados por una superficie reflectiva compuesta por múltiples facetas planas que reflejan los rayos solares horizontales y paralelos al eje óptico del concentrador y hacen seguimiento continuo del disco solar. El horno solar de la PSA consta de cuatro helióstatos dispuestos en dos niveles, cada uno de los cuales enfoca a una esquina del concentrador, de manera que se asegura la iluminación completa del concentrador durante el periodo operativo.

Helióstato
[26]


PS10: La megatorre sevillana de la energía solar (y cómo funciona)
 

Planta Solúcar PS10
En Sanlúcar la Mayor, a 18 kilómetros de Sevilla, la empresa española Abengoa construyó una estación solar de generación de electricidad.

Campo de helióstatos y torre
Inaugurada en el año 2007, la PS10 fue la primera central solar termoeléctrica (de carácter comercial) de torre central y campo de helióstatos instalada en el mundo. Utiliza paneles de espejo para generar 11 megavatios.
Sólo PS10, PS20 y Solnovas, que se reúnen en la conocida como Planta Solúcar, operan comercialmente un total de 183 MW, produciendo energía anual equivalente a 94.000 hogares, y evitando así la emisión de más de 114.000 toneladas anuales de CO2 a la atmósfera.
La PS10 está formada por 624 helióstatos y una torre solar de 114 metros de altura. El campo de espejos o helióstatos refleja la luz solar sobre un receptor en la parte superior de la torre.

Fascinante resplandor mágico
Esta torre se encuentra en el centro de la estación PS10 Solúcar. En la parte superior, el receptor solar consiste en una serie de paneles de tubos que operan a muy alta temperatura y por los que circula agua a presión. Este receptor se calienta por efecto de la luz solar y genera vapor saturado a 257 ºC. El vapor que se produce es almacenado parcialmente en unos tanques acumuladores para ser utilizado cuando no haya suficiente producción; el resto es enviado a accionar una turbina que genera la electricidad.

Esquema de funcionamiento
Los 624 helióstatos circundantes, cada uno con una superficie de 120 metros cuadrados, producen el reflejo para que el sistema pueda convertir alrededor del 17 por ciento de la energía de la luz solar en 11 megavatios de electricidad. Como punto de comparación, la planta PS2 (que se encuentra al lado) produce 20 MW de potencia mediante su torre de 160 metros de altura sobre un campo de 1.255 helióstatos.
Por lo tanto, esta planta funciona calentando agua con luz solar y, con el vapor generado, mueve turbinas para crear electricidad.

Campo y torre PS10
[14]



PS20
La PS20 (Planta Solar 20) es ahora mismo la mayor planta comercial de torre del mundo. Está situada al lado de la PS10.
A diferencia de la PS10, la PS20 cuenta con 85 hectáreas y con 1.255 helióstatos que reflejan los rayos al receptor situado en la parte superior de la torre de 160 metros de altura.
La PS20 dispone con un sistema de almacenamiento de 1 hora, con un receptor más eficiente y un sistema de control y operación mejor que le permiten producir 20 MW que abastecen a 10.000 hogares anualmente.

Helióstato PS10

[29]

Gemasolar
Gemasolar es una planta de energía solar por concentración con sistema de almacenamiento térmico en sales fundidas situada en Sevilla, España. Es la primera a escala comercial en aplicar la tecnología de receptor de torre central y almacenamiento térmico en sales fundidas. La superficie de esta planta ocupa aproximadamente 158 hectáreas. Características:
Potencia eléctrica de 19,9 MW.
La producción neta esperada es de 110 GWh/año.
Capaz de suministrar energía limpia y segura que reduce en más de 50.000 toneladas al año de emisiones de CO2.
La planta está formada por:
2.650 heliostatos que forman círculos concéntricos alrededor de la torre que disponen de un mecanismo que posiciona con precisión la superficie de los espejos.
Una torre que en lo alto se encuentra el receptor de haz de luz compuesto por paneles.
2 tanques: uno de sales frías y otro de sales calientes.
Un cambiador de calor; un generador y transformador.
Su funcionamiento se ilustra en la siguiente imagen:

Tanque 1: Sales frías.
Se bombean las sales a lo largo de la torre.
En el receptor de haz de luz, las sales se calientan y bajan al tanque 2 donde se almacenan a temperaturas superiores a 500º.
Tanque 2: Sales calientes.
Cambiador de agua.
Las sales al perder calor generan vapor de agua.
El vapor de agua hace que se mueva la turbina y el generador que produce la energía.
Transformador. La energía pasa al tendido eléctrico.
La planta no solo funciona como una central solar sino también como una central térmica (cuando no hay luz solar) gracias a su sistema de concentración de sales y a partir de ahí también generar electricidad. También podemos decir que gracias a esto la planta asegura obtención de energía las 24 horas del días en situaciones de baja insolación, en las madrugadas y varios meses al año.
Es tal el éxito de la planta sevillana que muchos profesionales del sector buscan aplicar las energías limpias a las pequeñas urbanizaciones y así hacer posible su aplicación a la vida cotidiana. [30]



Instalación de energía solar Ivanpah

El Sistema de generación eléctrica solar Ivanpah es una planta termosolar concentrada en el desierto de Mojave. Está ubicado en la base de Clark Mountain en California, a través de la línea estatal desde Primm, Nevada. La planta tiene una capacidad bruta de 392 megavatios (MW). Se despliega 173,500 helióstatos, cada uno con dos espejos de enfoque de la energía solar en calderas situados en tres centralizados torres de energía solar. La primera unidad del sistema se conectó a la red eléctrica en septiembre de 2013 para una prueba de sincronización inicial. La instalación abrió formalmente el 13 de febrero de 2014. En 2014, fue la estación de energía solar térmica más grande del mundo.
La instalación, con un costo de 2.2 mil millones de dólares, fue desarrollada por BrightSource Energy y Bechtel.
 
Descripción
El sistema Ivanpah consta de tres plantas de energía solar térmica en 4,000 acres (1,600 ha) de terrenos públicos cerca de la frontera entre California y Nevada en el suroeste de los Estados Unidos.
Los campos de espejos de helióstatos enfocan la luz solar en los receptores ubicados en torres de energía solar centralizadas. Los receptores generan vapor para accionar turbinas de vapor especialmente adaptadas.
Para la primera planta, se ordenó el mayor grupo generador de turbina de vapor con energía solar, con una turbina de recalentamiento de caja única Siemens SST-900 de 123 MW. Las plantas no tienen almacenaje, y calientan el vapor a 550 °C directamente en los receptores.
La aprobación final para el proyecto se otorgó en octubre de 2010. El 27 de octubre de 2010, el gobernador de California, Arnold Schwarzenegger, el secretario del Interior Ken Salazar y otros dignatarios se reunieron en el desierto de Mojave para abrir la tierra para la construcción.
 
Consumo de combustibles fósiles
La planta quema gas natural cada mañana para comenzar la operación. El 27 de agosto de 2014, el Estado de California aprobó a Ivanpah aumentar su consumo anual de gas natural de 328 a 525 millones de pies cúbicos. En 2014, la planta quemó 867,740 millones de BTU de gas natural que emitía 46,084 toneladas métricas de dióxido de carbono. En 2015, las instalaciones mostraron mayores números de producción, con aumentos del primer trimestre del 170 por ciento en el mismo período de 2014.
En 2015, el consumo de gas natural disminuyó a 564,814 millones de BTU, mientras que la producción total de energía aumentó a 652,300 MWh.
La instalación utiliza tres calderas de tubos de agua tipo D de Rentech y tres calderas de conservación durante la noche.
 
Impacto económico
BrightSource estimó que las instalaciones de Ivanpah proporcionarían 1,000 empleos en el pico de la construcción, 86 empleos permanentes y beneficios económicos totales de 3 mil millones de dólares.
El costo estimado de construcción para la instalación fue de 5,561.00 de dólares por kW.
En noviembre de 2014, los inversionistas de la instalación solicitaron una subvención federal de 539 millones de dólares para financiar su préstamo federal.
 
Rendimiento
El rendimiento de entrega de potencia contratada de 640 GW·h/año de las Unidades 1 y 3 y 336 GW·h de la Unidad 2 se cumplió en 2017, luego de una reducción drástica de la producción en los primeros años de operación, particularmente en el inicio del año 2014.
En noviembre de 2014, Associated Press informó que la instalación producía solo la mitad de su producción anual esperada. La Comisión de Energía de California emitió una declaración culpando a nubes, estelas de aviones y clima. El rendimiento mejoró en 2015 a aproximadamente 650 GW·h.
Para 2017, debido a mejoras, la planta estaba cumpliendo con los requisitos de producción del contrato.
 
Impactos ambientales
El proyecto generó controversia debido a la decisión de construirlo en un hábitat desértico ecológicamente intacto. La instalación de Ivanpah estimó reducir las emisiones de dióxido de carbono en más de 400,000 toneladas anuales. Fue diseñado para minimizar los impactos en el entorno natural. La instalación fue cercada para mantener alejada a la vida silvestre terrestre. Sin embargo, las aves se enfrentaron al riesgo de colisión con los espejos de helióstatos o por la combustión del flujo solar creado por el campo de espejos.
En 2012, la Asociación de Conservación de Parques Nacionales (NPCA) emitió un informe sobre el proyecto, citando preocupaciones sobre el agua, daños a los recursos visuales e impactos en importantes especies del desierto. Para conservar el escaso agua del desierto, el LPT 550 utiliza refrigeración por aire para convertir el vapor nuevamente en agua, y se devuelve a la caldera en un proceso cerrado.
Otro problema potencial es el efecto del reflejo del espejo en los pilotos de aviones. Las torres de energía tienen 'unidades receptoras' en su parte superior en las que los campos del espejo enfocan su luz reflejada. Durante la operación, estas unidades receptoras se vuelven extremadamente calientes, de modo que brillan y aparecen muy iluminadas. Debido a que están muy por encima del suelo, estas unidades receptoras brillantes son una distracción visible para las personas en muchos de los KOP (puntos de observación clave). [45]
 
 
Celóstato

En el siglo XIX, cuando no existían fuentes de luz artificial intensas como lámparas eléctricas o arcos voltaicos, era muy común utilizar la luz del sol en los experimentos de óptica. La forma más simple de obtener un haz de luz en un laboratorio consistía en colocar un espejo en el exterior, de manera que reflejara la luz del sol a la sala donde se realizaba el experimento. Si se dotaba al espejo de un movimiento de rotación adecuado se podía conseguir que la dirección del haz permaneciera constante durante varias horas, independientemente de la posición aparente del sol. En este caso, se tendría un helióstato.

Si se añade un segundo espejo, regulable en altura y orientación, se conseguía un celóstato. Este segundo espejo es importante porque la declinación del sol va variando a lo largo del año y por tanto la dirección del haz se proyectará a diferentes alturas. Con este segundo espejo se puede recoger la luz del primer espejo y enviarla a un instrumento astronómico, asegurando que la luz irá siempre al mismo sitio independientemente de la hora del día y de la época del año. [4]



RE<C: resumen del proyecto de helióstatos

Introducción

RE<C fue una iniciativa de Google para impulsar la innovación en energía renovable, con el objetivo de hacer que la energía renovable sea lo suficientemente barata para competir cara a cara con las centrales eléctricas de carbón. Google formó un equipo de ingeniería para desarrollar tecnologías prometedoras en el campo de la generación de energía solar. Centramos entonces nuestros esfuerzos de ingeniería en concentrar la energía solar (CSP).

Las plantas de energía solar concentradas usan espejos o lentes para enfocar una gran cantidad de luz solar sobre un objetivo que absorbe calor, llamado receptor. El intercambiador de calor del receptor crea vapor a alta presión, que luego impulsa una turbina para alimentar un generador eléctrico. La refrigeración por agua en spray se utiliza normalmente para condensar el vapor. Estas plantas de energía CSP suelen tener una potencia de 50 MW o más. El uso del agua puede ser un factor que limite la adopción de plantas de CSP a vapor, ya que el área con la mayor cantidad de luz solar que sería ideal para CSP a menudo tiene recursos hídricos limitados.

La CSP modular de "torre de energía" utiliza un motor de turbina de gas más pequeño (Brayton) para realizar la conversión de energía. Esta turbomáquina es similar en tamaño a los turbocompresores para motores diésel marinos o de camiones grandes. Los motores Brayton no necesitan refrigeración por rociado con agua y, de ese modo, se adaptan mejor a los ambientes secos del desierto.

El otro componente principal de una planta de energía CSP es un campo de espejos controlados, llamados helióstatos. Este campo tiene miles de metros cuadrados de helióstatos que concentran la energía solar en el receptor de la central eléctrica.

Nuestro diseño del helióstato

El módulo reflector

Cada helióstato tenía un espejo de enfoque de 2m x 3m articulado en la parte superior de su marco. Se utilizó un diseño liviano. El módulo reflector del espejo estaba hecho de vidrio. Se realizaron pruebas de granizo para verificar el cumplimiento de los estándares de la industria.

El marco del helióstato y la base

Muchos marcos y bases de helióstatos existentes son estructuras sólidas montadas sobre una base de hormigón vertido en un sitio plano. Utilizan accionamientos de precisión y grandes actuadores para realizar apuntamientos rígidos. Para reducir el costo, nuestro diseño tenía un marco liviano y fácil de transportar. Fue sujetado por un anclaje de tierra. Montados en el bastidor había dos accionadores de cable que usaban pequeños motores baratos. El helióstato tenía una junta en U, que se articulaba en inclinación/balanceo.

El diseño del campo.

Cada uno de nuestros sistemas modulares de conversión de potencia del motor Brayton fue diseñado para producir una salida eléctrica planificada de 890kW por torre, lo que requeriría 2600kW de energía solar térmica proveniente de la apertura del receptor.

Utilizamos nuestro software de simulación de helióstatos ópticos (HOpS) para experimentar con diferentes configuraciones de campo.

Establecimos un tamaño de campo de 862 helióstatos alrededor de una torre de 44 m, cada helióstato es de aproximadamente 6m2. Los helióstatos están dispuestos en un patrón hexagonal, a una distancia horizontal máxima de 60 m desde la torre.

Consideramos que los helióstatos tienen la funcionalidad de apuntar a múltiples torres.

Requisitos de focalización del sistema de control

Para convertir la energía de manera eficiente, el motor Brayton requiere un receptor de cavidad de temperatura más alta que el receptor típico de una planta de vapor. Se deben producir cambios pequeños e infrecuentes en el flujo de calor a través de su intercambiador de calor para prolongar su vida útil.

Un receptor de mayor temperatura requiere una abertura más pequeña para reducir la pérdida de calor radiante. Para prolongar la vida útil del receptor de alta temperatura, el flujo debe distribuirse cuidadosamente dentro de la cavidad del receptor.

Sistema de detección y control

El sistema de control es capaz de controlar simultáneamente los puntos de luz desde múltiples helióstatos a un alto grado de precisión a un lugar deseado en un objetivo. Fue rastreado para compensar el movimiento del sol a través del cielo mientras corrige los efectos del viento constante. No es sensible a los efectos de los cambios de cimentación o la expansión térmica del marco.

Se utilizó un acelerómetro de 3 ejes montado en helióstato de bajo costo combinado con un sistema central de fotometría multiscópica para resolver las posiciones de puntos de luz individuales en el objetivo.

Mitigación del viento

El viento presenta un desafío de diseño particularmente difícil, especialmente cuando se trata de diseñar helióstatos más ligeros y de bajo costo. Las áreas de tierra grandes y planas donde es más probable que se construyan los helióstatos son también las áreas más propensas al viento sin restricciones.

Los helióstatos a lo largo del borde exterior de un campo protegen a los helióstatos en el medio de gran parte del impacto del viento. Además, las cercas de viento simples pueden reducir dramáticamente el impacto del viento en los helióstatos. [43]



Tecnología solar de alta concentración

La actividad del Grupo de Alta Concentración Solar (GACS) se centra fundamentalmente en los sistemas de Receptor Central. El despliegue comercial de plantas solares termoeléctricas de Receptor Central (STE-RC) se inició tímidamente en España, con la inauguración de PS10 (2007) y PS20 (2009). A nivel internacional se observa, a finales de 2010, un renovado interés por las plantas de foco puntual.
La “curva de aprendizaje” de la solar termoeléctrica de Receptor central se basaba en el ensayo de más de 10 instalaciones experimentales de receptor central en el mundo y una amplia variedad de componentes.
La experiencia acumulada ha servido para demostrar la viabilidad técnica del concepto y su capacidad para trabajar a altas temperaturas e integrarse en ciclos más eficientes de forma escalonada. También han demostrado que admiten fácilmente el funcionamiento híbrido en varias opciones y tienen el potencial de generar electricidad mediante el uso de almacenamiento térmico.
Las plantas de Receptor Central ofrecen mayores eficiencias totales de conversión y una mayor diversidad de opciones de diseño, con menor experiencia acumulada en la implementación de cada tipología o componente. Pero tienen un elevado coste de inversión.
Este objetivo de reducir los costes de producción de electricidad está orientando los esfuerzos de I+D, para mejorar las opciones existentes y demostrar la viabilidad de nuevas opciones de diseño. [44]


\end{document}


